\documentclass[11pt]{article}
\usepackage[a4paper,margin=1in]{geometry}
\usepackage{fontspec}
\usepackage[boldfont=false]{xeCJK}
\setCJKmainfont{FandolSong-Regular.otf}
\usepackage{amsmath}
\usepackage{amssymb}
\usepackage{array}
\usepackage{booktabs}
\usepackage{graphicx}
\usepackage{adjustbox}
\usepackage{enumitem}
\setlist[itemize]{topsep=2pt,partopsep=0pt,parsep=0pt,itemsep=2pt,leftmargin=1.4em}
\setlist[enumerate]{topsep=2pt,partopsep=0pt,parsep=0pt,itemsep=2pt,leftmargin=1.6em}
\usepackage{parskip}
\usepackage{fvextra}
\fvset{breaklines=true,breakanywhere=true,fontsize=\small}
\setlength{\parindent}{0pt}

\begin{document}

\begin{center}
  {\LARGE\bfseries Forces, density and pressure}\\[0.35em]
  {\large A Level Physics}
\end{center}
\vspace{0.6em}

\section*{Turning effects of forces}

\subsection*{Centre of gravity}

The \textbf{weight} 重力 of a large object can be treated as acting at one single point, called the \textbf{centre of gravity} 重心 (the same as the \textbf{centre of mass} 质心 in a uniform gravitational field). For a uniform, regular shape — a rectangle, a sphere, a uniform rod — the centre of gravity is at the middle.

When you draw a \textbf{free-body diagram} 受力图, always put the weight arrow at the centre of gravity.

\subsection*{Moment of a force}

The \textbf{moment} 力矩 of a force about a point is

\[
M = F \cdot d,
\]

where $F$ is the size of the \textbf{force} 力 and $d$ is the \textbf{perpendicular distance} 垂直距离 from the point to the \textbf{line of action} 作用线 of the force. Unit: $\text{N m}$.

If the force acts at angle $\theta$ to a \textbf{lever arm} 力臂 of length $r$ from the \textbf{pivot} 支点, then $d = r\sin\theta$, so $M = F r\sin\theta$. Only the part of the force \textbf{perpendicular} 垂直 to the lever arm makes it turn.

\par\medskip
\noindent\begin{minipage}{\linewidth}
\centering
\adjustbox{max width=\linewidth, max totalheight=0.42\textheight}{%
  \includegraphics{/Users/joschka/Desktop/Programming/eXercise/output/eXam/handouts/a_level_physics/assets/04-p90-moment-of-a-force.pdf}}\\[0.3em]
\emph{The moment depends on the perpendicular distance d from the pivot to the line of action of the force}
\end{minipage}\par\medskip

A moment is either \textbf{clockwise} 顺时针 or \textbf{anticlockwise} 逆时针 about the chosen point.

\subsection*{Couple and torque}

A \textbf{couple} 力偶 is a pair of forces that are:

\begin{itemize}
\item \textbf{equal in size},

\item \textbf{opposite in direction},

\item with their lines of action a perpendicular distance apart.

\end{itemize}

A couple makes the body \textbf{turn only} — its resultant force is zero, so it gives no straight-line acceleration.

The \textbf{torque} 力偶矩 of a couple is the turning effect it makes:

\[
\tau = F \cdot d,
\]

where $F$ is the size of one force and $d$ is the perpendicular distance between the two lines of action. Unit: $\text{N m}$. The torque is the same about any point — a special property of couples.

\par\medskip
\noindent\begin{minipage}{\linewidth}
\centering
\adjustbox{max width=\linewidth, max totalheight=0.42\textheight}{%
  \includegraphics{/Users/joschka/Desktop/Programming/eXercise/output/eXam/handouts/a_level_physics/assets/04-p91-torque-of-a-couple.pdf}}\\[0.3em]
\emph{A couple: two equal and opposite forces, a distance apart, producing a torque}
\end{minipage}\par\medskip

A common multiple-choice trap: two equal forces in the \textbf{same} direction are not a couple (they have a resultant force and cause \textbf{translation} 平动). A couple needs equal size \textbf{and} opposite direction.

\section*{Equilibrium of forces}

\subsection*{Conditions for equilibrium}

A body is in \textbf{equilibrium} 平衡 when:

\begin{enumerate}
\item the \textbf{resultant force} 合力 is zero (no straight-line acceleration), AND

\item the resultant moment about any point is zero (no \textbf{angular acceleration} 角加速度).

\end{enumerate}

Both must hold. A body with no resultant force can still be turning; a body with no resultant moment can still be moving in a straight line.

\subsection*{Principle of moments}

For a body that is not turning, \textbf{the total clockwise moment about any point equals the total anticlockwise moment about the same point}. This is the \textbf{principle of moments} 力矩原理.

To solve a balance problem:

\begin{enumerate}
\item Choose a pivot — usually where an unknown force acts, so that force drops out (its distance is zero).

\item List every force and its perpendicular distance from the pivot.

\item Set $\sum M_{\text{clockwise}} = \sum M_{\text{anticlockwise}}$.

\item Use $\sum F = 0$ if you need a second equation.

\end{enumerate}

A ruler balanced on a pivot with masses on each side is solved this way. For a heavy uniform rod, remember to include its weight acting at its centre of gravity.

\par\medskip
\noindent\begin{minipage}{\linewidth}
\centering
\adjustbox{max width=\linewidth, max totalheight=0.42\textheight}{%
  \includegraphics{/Users/joschka/Desktop/Programming/eXercise/output/eXam/handouts/a_level_physics/assets/04-p92-principle-of-moments.pdf}}\\[0.3em]
\emph{Weights on a rule balanced at a pivot — used to test the principle of moments}
\end{minipage}\par\medskip

\subsection*{Vector triangle}

Three forces in the same plane that are in equilibrium can be drawn as a \textbf{closed vector triangle} 矢量三角形 — drawn tip-to-tail, the three arrows come back to the start. This is a drawing method instead of splitting into \textbf{components} 分量.

Use the \textbf{sine rule} 正弦定理 or the \textbf{cosine rule} 余弦定理 on the triangle to find unknown sizes or directions, or draw the triangle to scale on graph paper.

You can also split each force into \textbf{horizontal} 水平 and \textbf{vertical} 竖直 components and set $\sum F_{x} = 0$ and $\sum F_{y} = 0$.

\section*{Density}

\textbf{Density} 密度 is the mass per unit volume:

\[
\rho = \frac{m}{V}.
\]

Unit: $\text{kg m}^{-3}$ (or $\text{g cm}^{-3}$; $1\ \text{g cm}^{-3} = 1000\ \text{kg m}^{-3}$). Density is a \textbf{scalar} 标量.

Some useful densities to know:

\begin{itemize}
\item water: $1000\ \text{kg m}^{-3}$

\item air at room conditions: $\sim 1.2\ \text{kg m}^{-3}$

\item iron / steel: $\sim 7800\ \text{kg m}^{-3}$

\end{itemize}

\section*{Pressure}

\textbf{Pressure} 压强 is the force per unit area, where the force acts at right angles to the surface:

\[
p = \frac{F}{A}.
\]

Unit: $\text{Pa} = \text{N m}^{-2}$. Pressure is a scalar.

\subsection*{Hydrostatic pressure}

Take a column of \textbf{fluid} 流体 with density $\rho$, \textbf{cross-sectional area} 横截面积 $A$ and height $\Delta h$. Its weight is

\[
W = m g = (\rho \cdot A \cdot \Delta h) \cdot g.
\]

This weight presses down on the area $A$ at the bottom, so the extra pressure at the bottom compared with the top is

\[
\Delta p = \frac{W}{A} = \rho g \Delta h.
\]

This is the \textbf{hydrostatic pressure} 流体静压强 equation. It depends only on the density of the fluid and the \textbf{depth} 深度 — the shape of the container does not matter.

\par\medskip
\noindent\begin{minipage}{\linewidth}
\centering
\adjustbox{max width=\linewidth, max totalheight=0.42\textheight}{%
  \includegraphics{/Users/joschka/Desktop/Programming/eXercise/output/eXam/handouts/a_level_physics/assets/04-p94-hydrostatic-pressure-column.pdf}}\\[0.3em]
\emph{A column of liquid of area A: its weight sets the extra pressure at the depth below}
\end{minipage}\par\medskip

For a submarine at depth $h$ below the surface, the pressure from the water is $\rho_{\text{seawater}}\, g\, h$. For the total pressure, add the \textbf{atmospheric pressure} 大气压强 at the surface (about $1.0 \times 10^{5}\ \text{Pa}$).

\subsection*{Upthrust and Archimedes' principle}

When an object is \textbf{submerged} 浸没 in a fluid, the pressure at the bottom of the object is greater than the pressure at the top (by $\rho g \Delta h$, where $\Delta h$ is the object's height). This difference gives a net upward force called the \textbf{upthrust} 浮力.

\par\medskip
\noindent\begin{minipage}{\linewidth}
\centering
\adjustbox{max width=\linewidth, max totalheight=0.42\textheight}{%
  \includegraphics{/Users/joschka/Desktop/Programming/eXercise/output/eXam/handouts/a_level_physics/assets/04-p95-upthrust-origin.pdf}}\\[0.3em]
\emph{Upthrust arises because the pressure on the bottom of the object is greater than on the top}
\end{minipage}\par\medskip

For an object of volume $V$ (the volume of fluid \textbf{displaced} 排开), the upthrust is

\[
F_{\text{upthrust}} = \rho_{\text{fluid}}\, g\, V.
\]

This is \textbf{Archimedes' principle} 阿基米德原理: the upthrust on a body in a fluid equals the weight of the fluid it pushes aside.

For a fully submerged object, $V$ is its full volume. For a \textbf{floating} 漂浮 object, $V$ is only the volume below the surface — the object floats when the upthrust on the part below the surface equals its weight.

\subsection*{Force balance with upthrust}

A block held under water by a string tied to the bottom of the container is in equilibrium under three vertical forces: weight (down), \textbf{tension} 张力 (down), upthrust (up). Set $F_{\text{upthrust}} = W + T$ to find the tension.

A submerged block hanging from a \textbf{newton meter} 弹簧测力计 reads less than its weight in air, because of the upthrust: reading $= W - F_{\text{upthrust}}$.

The upthrust depends on the fluid density and the displaced volume, not on the object's material or depth (for an \textbf{incompressible} 不可压缩 fluid). On a planet with smaller $g$, the upthrust is smaller in the same ratio as the weight, so a floating object still floats with the same fraction below the surface.


\end{document}
